\documentclass[11pt]{amsart}
\usepackage[margin=1.1in]{geometry}
\usepackage{lmodern}
\usepackage{amsmath,amssymb,amsthm,mathtools}
\usepackage[mathscr]{eucal}
\usepackage{booktabs,enumitem,xcolor,array}
\usepackage[colorlinks=true,linkcolor=blue!50!black,citecolor=green!40!black,urlcolor=blue!60!black]{hyperref}
\hypersetup{
  pdftitle={Integer-weight solutions of the explicit formula: a characterisation of zeta among uniformly discrete Beurling systems},
  pdfauthor={Nic Johns}
}
\setlength{\emergencystretch}{2.5em}

% ---------------------------------------------------------------- macros
\newcommand{\RR}{\mathbb R}
\newcommand{\ZZ}{\mathbb Z}
\newcommand{\CC}{\mathbb C}
\newcommand{\NN}{\mathbb N}
\newcommand{\QQ}{\mathbb Q}
\newcommand{\dd}{\,\mathrm d}
\renewcommand{\Re}{\operatorname{Re}}
\renewcommand{\Im}{\operatorname{Im}}
\newcommand{\supp}{\operatorname{supp}}
\newcommand{\hF}{\widehat F}
\newcommand{\xin}[1]{\xi_{#1}}
\newcommand{\Zz}{Z_\zeta}
\newcommand{\muz}{\mu_\zeta}
\newcommand{\nuz}{\nu_\zeta}
\newcommand{\pz}{p_\zeta}
\newcommand{\Oinf}{\Omega_\infty}
\newcommand{\Arch}{\mathcal A}
\newcommand{\Tc}{\mathcal T}
\newcommand{\Spec}{\mathcal R}
\newcommand{\Kad}{\mathcal K}
\newcommand{\condU}{\textup{(U)}}
\newcommand{\RH}{\mathrm{RH}}

\newtheorem{theorem}{Theorem}[section]
\newtheorem{proposition}[theorem]{Proposition}
\newtheorem{lemma}[theorem]{Lemma}
\newtheorem{corollary}[theorem]{Corollary}
\newtheorem{fact}[theorem]{Fact}
\newtheorem{question}[theorem]{Question}
\theoremstyle{definition}
\newtheorem{definition}[theorem]{Definition}
\theoremstyle{remark}
\newtheorem{remark}[theorem]{Remark}
\numberwithin{equation}{section}

\title[Integer-weight solutions of the explicit formula]{Integer-weight solutions of the explicit formula:\\
a characterisation of $\zeta$ among uniformly discrete Beurling systems}
\author{Nic Johns}
\date{October 2026}
\subjclass[2020]{Primary 11M26; Secondary 11M41, 11N80, 42A38}
\keywords{Weil explicit formula, Hamburger's theorem, Beurling generalised integers, Fourier quasicrystals}

\begin{document}

\begin{abstract}
A positive solution of Weil's explicit formula with the archimedean data of the Riemann zeta function is a pair of
positive measures, one on the zero side and one on the prime side beyond $\log2/2\pi$, that satisfies the formula on a
natural test class; the pair formed by the zeros and the prime powers of $\zeta$ is such a solution exactly when the
Riemann hypothesis holds. We consider solutions whose zero measure has integer weights, as for an $L$-function. If the
prime measure is carried by the logarithms of the integers, Hamburger's theorem shows that the solution is $\zeta$'s
own, and the Riemann hypothesis follows. We prove the same conclusion when the prime measure generates a uniformly
discrete multiplicative monoid, that is, a uniformly discrete Beurling system with no generalised primes below $2$.
Such a solution comes from a generalised Dirichlet series with Riemann's functional equation; the theorem of Lev and
Olevskii on measures with uniformly discrete support and spectrum, combined with an argument on linear recurrences in
a number field, forces the monoid into the positive integers.
\end{abstract}

\maketitle

\noindent\textbf{Status.} This note accompanies the companion paper \cite{PaperI} and has not been peer-reviewed. Its
theorem is proved in full, and no computation is used. Part (a) is Hamburger's theorem in the language of \cite{PaperI};
to our knowledge part (b) is new. The definitions and the results of \cite{PaperI} that are used are recalled in
Section~\ref{sec:setting}, as Facts.

%======================================================================
\section{Introduction}\label{sec:intro}
%======================================================================

An \emph{admissible pair} \cite{PaperI} is a pair $(\mu,\nu)$ of positive measures, $\mu$ even on $\RR$ and $\nu$ on
$[\xin2,\infty)$ with $\xin2=\log2/2\pi$, that satisfies Weil's explicit formula with the archimedean data of $\zeta$ on a
natural test class (Section~\ref{sec:setting}). The pair $\pz$ of $\zeta$ itself, formed by the real zero ordinates with
their multiplicities and by the measure $\sum_n\Lambda(n)n^{-1/2}\delta_{\xin n}$, is admissible if and only if the Riemann
hypothesis holds, and \cite{PaperI} conjectures that every admissible pair, if one exists, is this pair (Conjecture~U). An admissible zero
measure need not be atomic, and its atoms need not be integers. This note treats the case in which $\mu$ has integer
weights, as the zero measure of an $L$-function has. For a positive measure $\nu$ on $[\xin2,\infty)$ let $\tilde\nu$ be its
image under $\xi\mapsto e^{2\pi\xi}$, a measure on $[2,\infty)$.

\begin{theorem}\label{thm:integer}
Let $(\mu,\nu)$ be an admissible pair, and suppose that $\mu$ is purely atomic with weights in $\ZZ_{\ge0}$.
\begin{enumerate}[label=\textup{(\alph*)}]
\item \textup{(Hamburger)} If $\nu$ is carried by $\{\xin n:n\ge2\text{ an integer}\}$, then the Riemann hypothesis holds and
$(\mu,\nu)=\pz$.
\item If the multiplicative monoid $\mathcal N\subset[1,\infty)$ generated by $\{e^{2\pi\xi}:\xi\in\supp\nu\}$ \textup{(}the set of finite
products of these numbers, including the empty product $1$\textup{)} is uniformly
discrete, then the Riemann hypothesis holds and $(\mu,\nu)=\pz$.
\end{enumerate}
\end{theorem}

The proof is in Section~\ref{sec:proof}. Testing the explicit formula on Poisson kernels shows that
$L(s):=\int y^{1/2-s}\dd\tilde\nu(y)$ is the negative logarithmic derivative of a generalised Dirichlet series with positive
coefficients, Riemann's functional equation and finite order. Part (a) is then Hamburger's theorem. In (b), the theorem of
Lev and Olevskii on measures with uniformly discrete support and spectrum puts the monoid into finitely many cosets of
a lattice, and an argument with linear recurrences in a real number field forces it into $\NN$, which reduces (b) to (a).

Part (a) is Hamburger's theorem in the language of admissible pairs; positivity of $\nu$ is used only for convergence. The
theorem of Chandrasekharan and Mandelbrojt \cite[Theorem~3]{CM57}, generalised by Kahane and Mandelbrojt
\cite[Theorem~7]{KM58}, says that the frequencies of a general Dirichlet series with Riemann's functional equation are
finitely generated; it assumes a gap condition on the frequencies, which \cite{KM58} weakens to conditions of density type. This is too weak for the rationality step in (b), since the powers of $1+\sqrt2$ lie in $\ZZ[\sqrt2]$; the
lattice structure supplied by Lev and Olevskii is what is needed. Bochner and Chandrasekharan \cite{BC56} proved uniqueness
for general Dirichlet series with Riemann's functional equation under density restrictions on the frequencies. Beurling
\cite{Beu51} characterised $\zeta(s)$ and $(2^s-1)\zeta(s)$ among general Dirichlet series $\sum\lambda_n^{-s}$ with unit
coefficients by the pole at $s=1$, the trivial zeros and a growth condition. A.~B.~Dixit \cite{DixA20} showed that an element
$F$ of the extended Selberg class with non-negative coefficients and a simple pole at $s=1$ is determined, up to at most $2d_F$
possibilities, by its degree $d_F$, conductor and residue, under an explicit inequality between the conductor and the residue.
Hilberdink \cite{Hil12} proved rigidity results for Beurling generalised integers. Lagarias \cite{Lag99} studied Beurling
generalised integers with the Delone property, $r\le n_{i+1}-n_i\le R$, and classified those contained in $\ZZ$: their generalised
primes are all but finitely many primes, together with finitely many composites. Theorem~\ref{thm:integer}(b) assumes only the
lower bound, and not that $\mathcal N\subset\ZZ$: Step~5 derives $\mathcal N\subset\NN$ from the functional equation, and Hamburger's
theorem then leaves only $\mathcal N=\NN$, so among the semigroups classified by Lagarias the explicit formula with $\zeta$'s data
singles out $\NN$ itself. Related classifications, of Fourier quasicrystals with unit masses \cite{OU20} and of summation formulas
whose measure has uniformly discrete support and masses bounded below \cite[Theorem~5]{Gon23}, describe the support as the zero set
of an entire function of a special form. These results impose positivity or
integrality on the Dirichlet coefficients; Theorem~\ref{thm:integer} imposes integrality on the zero measure and works with the
explicit formula. With integer atoms, the uniqueness of the admissible pair becomes a uniqueness problem for Beurling
generalised number systems \cite{Beu37,DZ16} with no generalised primes below $2$ and Riemann's functional equation, and part
(b) settles it for the systems whose integers are uniformly discrete. The general case is open (Section~\ref{sec:open}).

%======================================================================
\section{Setting}\label{sec:setting}
%======================================================================

We recall the definitions and results of \cite{PaperI} that are used; statements quoted from it are called Facts. We write
$\hF(\xi)=\int_\RR F(t)e^{-2\pi i\xi t}\dd t$ and $\xin n=\log n/(2\pi)$ for real $n>0$. The archimedean data of $\zeta$ are
\[
\Oinf(t):=\frac1{2\pi}\Big(\Re\psi\big(\tfrac14+\tfrac{it}2\big)-\log\pi\Big),\qquad
\Arch(F):=F\big(\tfrac i2\big)+F\big(-\tfrac i2\big)+\int_\RR F(t)\Oinf(t)\dd t,
\]
with $\psi=\Gamma'/\Gamma$. For $0<\delta<\frac12$, $\Tc_\delta$ is the set of even functions $F$, real on $\RR$ and analytic in a
neighbourhood of $\{\lvert\Im z\rvert\le\frac12+\delta\}$, with $\sup_{\lvert\Im z\rvert\le\frac12+\delta}(1+\lvert z\rvert)^2\lvert F(z)\rvert<\infty$, and
$\Tc:=\bigcup_{\delta>0}\Tc_\delta$ is the \emph{test class}. It contains the Poisson kernels $t\mapsto(s-\frac12)/((s-\frac12)^2+t^2)$
for real $s>1$.

\begin{fact}[Explicit formula on $\Tc$ {\cite{PaperI}}; {\cite{Wei52}}, {\cite[(1.1)]{BRS}}]\label{lem:explicit-formula}
For every $F\in\Tc$, $\Arch(F)=\sum_\rho F(t_\rho)+\frac1\pi\sum_{n\ge2}\Lambda(n)n^{-1/2}\hF(\xin n)$, where $t_\rho:=-i(\rho-\frac12)$ and
$\rho$ runs over the non-trivial zeros of $\zeta$ with multiplicity. Both series converge absolutely; no hypothesis on the
zeros is made.
\end{fact}

Let $\Zz$ be the set of real $\gamma$ with $\zeta(\frac12+i\gamma)=0$, $m_\gamma$ the multiplicity,
$\muz:=\sum_{\gamma\in\Zz}m_\gamma\delta_\gamma$, $\nuz:=\sum_{n\ge2}\Lambda(n)n^{-1/2}\delta_{\xin n}$ and $\pz:=(\muz,\nuz)$.

\begin{definition}[Admissible pairs {\cite{PaperI}}]\label{def:admissible}
An \emph{admissible pair} is a pair $(\mu,\nu)$ of positive Radon measures, $\mu$ even on $\RR$ and $\nu$ on $[\xin2,\infty)$, such
that for every $F\in\Tc$ the integrals $\int F\dd\mu$ and $\int\hF\dd\nu$ converge absolutely and
\[
\Arch(F)=\Spec_{\mu,\nu}(F):=\int_\RR F\dd\mu+\frac1\pi\int_{[\xin2,\infty)}\hF\dd\nu .
\]
$\Kad$ denotes the set of admissible pairs.
\end{definition}

Under the Riemann hypothesis $\pz\in\Kad$, by Fact~\ref{lem:explicit-formula}. Conversely, if $(\mu,\nuz)\in\Kad$ for some $\mu$, then
the Riemann hypothesis holds and $\mu=\muz$ \cite[Theorem~2.9(b)]{PaperI}. So $\pz\in\Kad$ if and only if the Riemann hypothesis holds.

\begin{fact}[A priori bounds {\cite{PaperI}}]\label{lem:apriori}
There is an absolute constant $C$ such that every $(\mu,\nu)\in\Kad$ satisfies
\[
\mu([T-1,T+1])\le C\log(3+\lvert T\rvert)\quad(T\in\RR),\qquad \nu([X-1,X+1])\le Ce^{\pi X}\quad(X\in\RR).
\]
Hence $\int(1+t^2)^{-1}\dd\mu<\infty$ and $\int e^{-\pi(1+\eta)\xi}\dd\nu(\xi)<\infty$ for every $\eta>0$.
\end{fact}

%======================================================================
\section{Proof of Theorem~\ref{thm:integer}}\label{sec:proof}
%======================================================================

\begin{proof}
Write $\mu=\sum_\gamma w_\gamma\delta_\gamma$ with $w_\gamma\in\ZZ_{\ge0}$ and $w_{-\gamma}=w_\gamma$.

\emph{Step 1: the logarithmic derivative.} For real $s>1$ the Poisson kernel
$P_s(t):=\frac1\pi\frac{s-1/2}{(s-1/2)^2+t^2}$ lies in $\Tc$, with $\widehat{P_s}(\xi)=e^{-2\pi(s-1/2)|\xi|}$. We have
$P_s(\frac i2)+P_s(-\frac i2)=\frac1\pi(\frac1s+\frac1{s-1})$, and $\int P_s\Oinf=\frac1{2\pi}(\psi(\frac s2)-\log\pi)$: the function
$t\mapsto\psi(\frac14-\frac{it}2)$ is analytic with logarithmic growth on the closed upper half-plane, so its Poisson integral at
$t=i(s-\frac12)$ is its value there, $\psi(\frac s2)$, and similarly for $\psi(\frac14+\frac{it}2)$ in the lower half-plane. On the
spectral side $\int P_s\dd\mu=\frac1\pi\sum_\gamma w_\gamma(s-\frac12-i\gamma)^{-1}$ (summed symmetrically), and
$\frac1\pi\int\widehat{P_s}\dd\nu=\frac1\pi L(s)$ with $L(s):=\int y^{1/2-s}\dd\tilde\nu(y)$, which converges absolutely for $\Re s>1$ by
Fact~\ref{lem:apriori}. By analytic continuation, for $\Re s>1$,
\begin{equation}\label{eq:Ls}
L(s)=\frac1s+\frac1{s-1}+\frac12\psi\big(\tfrac s2\big)-\frac12\log\pi-\sum_\gamma\frac{w_\gamma}{s-\frac12-i\gamma}.
\end{equation}
The right side is meromorphic on $\CC$ (as $\sum_\gamma w_\gamma(1+\gamma^2)^{-1}<\infty$), with simple poles only: residue $1$ at
$s=1$, $0$ at $s=0$, $-1$ at $s=-2k$ $(k\ge1)$, and $-w_\gamma$ at $\frac12+i\gamma$. All residues are integers.

\emph{Step 2: a generalised Dirichlet series.} Let $\Pi:=y^{1/2}(\log y)^{-1}\tilde\nu$, a positive measure on $[2,\infty)$, and
$\mathcal Z(s):=\exp(\int y^{-s}\dd\Pi(y))$ for $\Re s>1$, so that $-\mathcal Z'/\mathcal Z=L$. Then $\mathcal Z(s)=\int y^{-s}\dd\mathsf N(y)$
with $\mathsf N:=\sum_{k\ge0}\Pi^{*k}/k!$ (multiplicative convolution, $\Pi^{*0}=\delta_1$), a positive measure on $[1,\infty)$ with
$\mathsf N(\{1\})=1$ whose support is the closure of the monoid generated by $\supp\tilde\nu$. Because the residues of $L$
are integers, $\mathcal Z$ continues to a meromorphic function on $\CC$ with a simple pole at $s=1$, zeros of order $w_\gamma$ at
$\frac12+i\gamma$, simple zeros at $s=-2k$, and no other zeros or poles.

\emph{Step 3: functional equation and order.} Put $\Lambda_{\mathcal Z}(s):=\pi^{-s/2}\Gamma(\frac s2)\mathcal Z(s)$. By \eqref{eq:Ls},
$R(s):=-\Lambda_{\mathcal Z}'/\Lambda_{\mathcal Z}(s)=\frac1s+\frac1{s-1}-\sum_\gamma w_\gamma(s-\frac12-i\gamma)^{-1}$, and $R(1-s)=-R(s)$
because $w$ is even. Hence $\Lambda_{\mathcal Z}(1-s)=C\Lambda_{\mathcal Z}(s)$ with a constant $C$, and $C^2=1$. The function
$E(s):=s(s-1)\Lambda_{\mathcal Z}(s)$ is entire with zeros exactly at the points $\rho_\gamma:=\frac12+i\gamma$, of orders $w_\gamma$. By
Fact~\ref{lem:apriori}, $\sum_{|\gamma|\le T}w_\gamma=O(T\log T)$, so the genus-one canonical product $P$ over these zeros has order at
most $1$; and $E'/E-P'/P=-\sum_\gamma w_\gamma/\rho_\gamma$ (summed symmetrically) is constant. So $E=e^{a+bs}P$ has order at most $1$,
and $(s-1)\mathcal Z(s)=E(s)\pi^{s/2}/(2\Gamma(1+\frac s2))$ is entire of finite order.

\emph{Step 4: proof of (a).} Now $\mathsf N$ is carried by the positive integers: $\mathcal Z(s)=\sum_{m\ge1}c_mm^{-s}$ with $c_1=1$,
absolutely convergent for $\Re s>1$. Hamburger's theorem in its two-series form \cite[Satz~1]{Ham21} (see also \cite{Bur12};
\cite[Theorem~2.1]{Per17} is the one-series form) applies to $f=\mathcal Z$ and $g=C\mathcal Z$: a Dirichlet series $f$ absolutely
convergent for $\Re s>1$, with $(s-1)f$ entire of finite order and $\pi^{-s/2}\Gamma(\frac s2)f(s)=\pi^{-(1-s)/2}\Gamma(\frac{1-s}2)g(1-s)$ for
a Dirichlet series $g$ absolutely convergent for $\Re s>1$, satisfies $f=g=c\zeta$. Since $c_1=1$, $\mathcal Z=\zeta$ and $C=1$. Then
$L=-\zeta'/\zeta$, so $\nu=\nuz$, and the zeros of $\zeta$ in $0<\Re s<1$ are the points $\frac12+i\gamma$ with multiplicities
$w_\gamma$: the Riemann hypothesis holds and $\mu=\muz$.

\emph{Step 5: proof of (b).} Here $\mathsf N$ is carried by the closed set $\mathcal N$, and by positivity its support is all of
$\mathcal N$: $\mathcal Z(s)=\sum_{n\in\mathcal N}c_nn^{-s}$ with $c_n>0$ and $c_1=1$. Let $\varrho>0$ be the residue of $\mathcal Z$ at $s=1$
(positive, since $\mathcal Z(\sigma)>0$ for $\sigma>1$). The functional equation of Step~3 and the analytic properties established
there are equivalent, by Mellin inversion, to the summation formula $\Theta(1/x)=C\sqrt x\,\Theta(x)$ $(x>0)$ for
$\Theta(x):=C\varrho+2\sum_{n\in\mathcal N}c_ne^{-\pi n^2x}$ \cite[Introduction, Theorems~1--3 and Propositions~3--4]{KM58}.
Since the Fourier transform of $e^{-\pi xt^2}$ is $x^{-1/2}e^{-\pi\xi^2/x}$, and the Gaussians $e^{-\pi xt^2}$ $(x>0)$ span a dense
subspace of the even Schwartz functions (the derivatives in $x$ at $x=1$, limits of difference quotients in the Schwartz topology,
are the functions $t^{2k}e^{-\pi t^2}$, whose span contains the even Hermite functions), the even tempered measure $m:=C\varrho\delta_0+\sum_{n\in\mathcal N}c_n(\delta_n+\delta_{-n})$
satisfies $\widehat m=Cm$. Its support $\{0\}\cup\pm\mathcal N$ is uniformly discrete, and so is its spectrum. By the theorem of Lev
and Olevskii \cite[Theorem~1]{LO15}, which holds for complex measures on $\RR$ (so the sign $C=-1$ is allowed), $\supp m$ lies in a
finite union of cosets $\lambda_j+\alpha\ZZ$.

The set $\mathcal N$ is infinite, since $\nu\ne0$ (otherwise $L\equiv0$). Hence some coset contains two elements of $\mathcal N$, and we
may assume that every coset meets $\mathcal N$ and that $\lambda_j\in\mathcal N$. The $\QQ$-span $K$ of $\mathcal N$ is a subring of $\RR$
(as $\mathcal N$ is multiplicative and $1\in\mathcal N$) of finite dimension over $\QQ$, so $K$ is a real number field, and $\alpha\in K$.
Fix $n\in\mathcal N$ with $n>1$ and a $\QQ$-linear $\ell:K\to\QQ$ vanishing on $\QQ\alpha$. Then $u_k:=\ell(n^k)$ takes only the finitely
many values $\ell(\lambda_j)$. Writing $\ell(x)=\mathrm{Tr}_{K/\QQ}(\beta x)$ and $b:=\mathrm{Tr}_{K/\QQ(n)}(\beta)$, we get
$u_k=\sum_\sigma\sigma(b)\sigma(n)^k$, the sum over the embeddings $\sigma$ of $\QQ(n)$. A linear recurrence sequence with finitely many
values is eventually periodic, say $u_{k+p}=u_k$ for $k\ge k_0$; as the $\sigma(n)$ are distinct and non-zero, Vandermonde's
determinant gives $\sigma(b)(\sigma(n)^p-1)=0$ for every $\sigma$. For the identity embedding $n^p\ne1$, so $b=0$ and $u\equiv0$. This holds
for every $\QQ$-linear $\ell$ vanishing on $\QQ\alpha$, and these forms separate the points of $K$ modulo $\QQ\alpha$. Hence
$n^k\in\QQ\alpha$ for all $k\ge0$; $k=0$ gives $\alpha\in\QQ$ and $k=1$ gives $n\in\QQ$. Thus $\mathcal N\subset\QQ$ and all $\lambda_j\in\QQ$,
so $\mathcal N\subset D^{-1}\ZZ$ for some $D\in\NN$. If $a/b\in\mathcal N$ in lowest terms had $b>1$, its powers would have unbounded
denominators. So $\mathcal N\subset\NN$, and (a) applies.
\end{proof}

%======================================================================
\section{Remarks and open problems}\label{sec:open}
%======================================================================

\begin{remark}[Atomicity at criticality]
The slack of \cite{PaperI} is $\kappa^*:=\inf\{\Arch(F):F\in\Tc,\ F\ge0\text{ on }\RR,\ \hF\ge0\text{ on }[\xin2,\infty),\ \int F=1\}$;
admissible pairs exist if and only if $\kappa^*\ge0$ \cite[Theorem~A]{PaperI}. If $\kappa^*=0$, every admissible zero measure is purely atomic
\cite{PaperI}, but the weights of its atoms need not be integers; Theorem~\ref{thm:integer} needs integer weights.
\end{remark}

\begin{enumerate}[label=\textup{(\arabic*)}]
\item\label{op:Beurling} \emph{The Beurling case.} Uniqueness of the admissible pair when $\mu$ has integer weights and the prime
measure generates a monoid that is not uniformly discrete; Theorem~\ref{thm:integer} leaves exactly this case open.
\item \emph{Weaker discreteness.} Is uniform discreteness of the generated monoid essential? Would discreteness of $\supp\tilde\nu$,
together with a polynomial bound on the counting function of the monoid, suffice, for instance through the theorem of Kahane
and Mandelbrojt \cite[Theorem~7]{KM58} and a uniqueness result for measures whose support is not uniformly discrete?
\end{enumerate}

\section*{AI-use statement}
This note was produced with substantial use of AI. The mathematical arguments and the text were generated by instances of Claude Opus 5.5
(Anthropic) working as coordinated agents: one set of agents developed the arguments, separate agents acted as
independent hostile referees of each component, and the referees' requested repairs were incorporated. The
author initiated and directed this project, designed its verification process, and reviewed the results and the
verification record, but did not independently check every step of the written proofs. The author takes full
responsibility for the content of this note, including any errors, and welcomes independent verification and
corrections.

\bibliographystyle{amsplain-doi}
\bibliography{refs}

\end{document}
